%DIF 1c1
%DIF LATEXDIFF DIFFERENCE FILE
%DIF DEL ../prl_submission/prl_main.tex   Wed Jul 22 09:14:24 2026
%DIF ADD prm_main.tex                     Sat Aug  8 11:13:07 2026
%DIF < \documentclass[aps,prl,twocolumn,10pt,superscriptaddress,floatfix]{revtex4-2}
%DIF -------
\documentclass[aps,prmaterials,twocolumn,10pt,superscriptaddress,floatfix]{revtex4-2} %DIF > 
%DIF -------

\usepackage{amsmath}
\usepackage{mathtools}
\usepackage{graphicx}
\usepackage{amsfonts}
\usepackage[colorlinks=true,allcolors=blue]{hyperref}
\Urlmuskip=0mu plus 2mu\relax

\graphicspath{{figures/}}
 %DIF > 
% Keep natural paragraph spacing; pool column slack at the bottom instead %DIF > 
% of stretching the parskip glue (visible in equation-free columns). %DIF > 
\raggedbottom %DIF > 
 %DIF > 
% Run-in section heads. Adds a little vertical space above the paragraph %DIF > 
% to separate the Letter's sections; the glue is discardable at column %DIF > 
% tops. Tune the amount here. %DIF > 
\newcommand{\runinsec}[1]{\medskip\emph{#1}} %DIF > 
%DIF PREAMBLE EXTENSION ADDED BY LATEXDIFF
%DIF UNDERLINE PREAMBLE %DIF PREAMBLE
\RequirePackage[normalem]{ulem} %DIF PREAMBLE
\RequirePackage{color}\definecolor{RED}{rgb}{1,0,0}\definecolor{BLUE}{rgb}{0,0,1} %DIF PREAMBLE
\providecommand{\DIFaddtex}[1]{{\protect\color{blue}\uwave{#1}}} %DIF PREAMBLE
\providecommand{\DIFdeltex}[1]{{\protect\color{red}\sout{#1}}}                      %DIF PREAMBLE
%DIF SAFE PREAMBLE %DIF PREAMBLE
\providecommand{\DIFaddbegin}{} %DIF PREAMBLE
\providecommand{\DIFaddend}{} %DIF PREAMBLE
\providecommand{\DIFdelbegin}{} %DIF PREAMBLE
\providecommand{\DIFdelend}{} %DIF PREAMBLE
\providecommand{\DIFmodbegin}{} %DIF PREAMBLE
\providecommand{\DIFmodend}{} %DIF PREAMBLE
%DIF FLOATSAFE PREAMBLE %DIF PREAMBLE
\providecommand{\DIFaddFL}[1]{\DIFadd{#1}} %DIF PREAMBLE
\providecommand{\DIFdelFL}[1]{\DIFdel{#1}} %DIF PREAMBLE
\providecommand{\DIFaddbeginFL}{} %DIF PREAMBLE
\providecommand{\DIFaddendFL}{} %DIF PREAMBLE
\providecommand{\DIFdelbeginFL}{} %DIF PREAMBLE
\providecommand{\DIFdelendFL}{} %DIF PREAMBLE
%DIF HYPERREF PREAMBLE %DIF PREAMBLE
\providecommand{\DIFadd}[1]{\texorpdfstring{\DIFaddtex{#1}}{#1}} %DIF PREAMBLE
\providecommand{\DIFdel}[1]{\texorpdfstring{\DIFdeltex{#1}}{}} %DIF PREAMBLE
\newcommand{\DIFscaledelfig}{0.5}
%DIF HIGHLIGHTGRAPHICS PREAMBLE %DIF PREAMBLE
\RequirePackage{settobox} %DIF PREAMBLE
\RequirePackage{letltxmacro} %DIF PREAMBLE
\newsavebox{\DIFdelgraphicsbox} %DIF PREAMBLE
\newlength{\DIFdelgraphicswidth} %DIF PREAMBLE
\newlength{\DIFdelgraphicsheight} %DIF PREAMBLE
% store original definition of \includegraphics %DIF PREAMBLE
\LetLtxMacro{\DIFOincludegraphics}{\includegraphics} %DIF PREAMBLE
\newcommand{\DIFaddincludegraphics}[2][]{{\color{blue}\fbox{\DIFOincludegraphics[#1]{#2}}}} %DIF PREAMBLE
\newcommand{\DIFdelincludegraphics}[2][]{% %DIF PREAMBLE
\sbox{\DIFdelgraphicsbox}{\DIFOincludegraphics[#1]{#2}}% %DIF PREAMBLE
\settoboxwidth{\DIFdelgraphicswidth}{\DIFdelgraphicsbox} %DIF PREAMBLE
\settoboxtotalheight{\DIFdelgraphicsheight}{\DIFdelgraphicsbox} %DIF PREAMBLE
\scalebox{\DIFscaledelfig}{% %DIF PREAMBLE
\parbox[b]{\DIFdelgraphicswidth}{\usebox{\DIFdelgraphicsbox}\\[-\baselineskip] \rule{\DIFdelgraphicswidth}{0em}}\llap{\resizebox{\DIFdelgraphicswidth}{\DIFdelgraphicsheight}{% %DIF PREAMBLE
\setlength{\unitlength}{\DIFdelgraphicswidth}% %DIF PREAMBLE
\begin{picture}(1,1)% %DIF PREAMBLE
\thicklines\linethickness{2pt} %DIF PREAMBLE
{\color[rgb]{1,0,0}\put(0,0){\framebox(1,1){}}}% %DIF PREAMBLE
{\color[rgb]{1,0,0}\put(0,0){\line( 1,1){1}}}% %DIF PREAMBLE
{\color[rgb]{1,0,0}\put(0,1){\line(1,-1){1}}}% %DIF PREAMBLE
\end{picture}% %DIF PREAMBLE
}\hspace*{3pt}}} %DIF PREAMBLE
} %DIF PREAMBLE
\LetLtxMacro{\DIFOaddbegin}{\DIFaddbegin} %DIF PREAMBLE
\LetLtxMacro{\DIFOaddend}{\DIFaddend} %DIF PREAMBLE
\LetLtxMacro{\DIFOdelbegin}{\DIFdelbegin} %DIF PREAMBLE
\LetLtxMacro{\DIFOdelend}{\DIFdelend} %DIF PREAMBLE
\DeclareRobustCommand{\DIFaddbegin}{\DIFOaddbegin \let\includegraphics\DIFaddincludegraphics} %DIF PREAMBLE
\DeclareRobustCommand{\DIFaddend}{\DIFOaddend \let\includegraphics\DIFOincludegraphics} %DIF PREAMBLE
\DeclareRobustCommand{\DIFdelbegin}{\DIFOdelbegin \let\includegraphics\DIFdelincludegraphics} %DIF PREAMBLE
\DeclareRobustCommand{\DIFdelend}{\DIFOaddend \let\includegraphics\DIFOincludegraphics} %DIF PREAMBLE
\LetLtxMacro{\DIFOaddbeginFL}{\DIFaddbeginFL} %DIF PREAMBLE
\LetLtxMacro{\DIFOaddendFL}{\DIFaddendFL} %DIF PREAMBLE
\LetLtxMacro{\DIFOdelbeginFL}{\DIFdelbeginFL} %DIF PREAMBLE
\LetLtxMacro{\DIFOdelendFL}{\DIFdelendFL} %DIF PREAMBLE
\DeclareRobustCommand{\DIFaddbeginFL}{\DIFOaddbeginFL \let\includegraphics\DIFaddincludegraphics} %DIF PREAMBLE
\DeclareRobustCommand{\DIFaddendFL}{\DIFOaddendFL \let\includegraphics\DIFOincludegraphics} %DIF PREAMBLE
\DeclareRobustCommand{\DIFdelbeginFL}{\DIFOdelbeginFL \let\includegraphics\DIFdelincludegraphics} %DIF PREAMBLE
\DeclareRobustCommand{\DIFdelendFL}{\DIFOaddendFL \let\includegraphics\DIFOincludegraphics} %DIF PREAMBLE
%DIF COLORLISTINGS PREAMBLE %DIF PREAMBLE
\RequirePackage{listings} %DIF PREAMBLE
\RequirePackage{color} %DIF PREAMBLE
\lstdefinelanguage{DIFcode}{ %DIF PREAMBLE
%DIF DIFCODE_UNDERLINE %DIF PREAMBLE
  moredelim=[il][\color{red}\sout]{\%DIF\ <\ }, %DIF PREAMBLE
  moredelim=[il][\color{blue}\uwave]{\%DIF\ >\ } %DIF PREAMBLE
} %DIF PREAMBLE
\lstdefinestyle{DIFverbatimstyle}{ %DIF PREAMBLE
	language=DIFcode, %DIF PREAMBLE
	basicstyle=\ttfamily, %DIF PREAMBLE
	columns=fullflexible, %DIF PREAMBLE
	keepspaces=true %DIF PREAMBLE
} %DIF PREAMBLE
\lstnewenvironment{DIFverbatim}{\lstset{style=DIFverbatimstyle}}{} %DIF PREAMBLE
\lstnewenvironment{DIFverbatim*}{\lstset{style=DIFverbatimstyle,showspaces=true}}{} %DIF PREAMBLE
%DIF END PREAMBLE EXTENSION ADDED BY LATEXDIFF

\begin{document}

\title{Yukawa screening derivation of the bond-valence rule}

\author{Michael L. Whittaker}
\email{mwhittaker@lbl.gov}
\affiliation{Energy Geosciences Division, Lawrence Berkeley National Laboratory, Berkeley, CA 94720}
\affiliation{Materials Sciences Division, Lawrence Berkeley National Laboratory, Berkeley, CA 94720}

\author{Pan Wang}
\affiliation{Energy Geosciences Division, Lawrence Berkeley National Laboratory, Berkeley, CA 94720}

\author{Chunhui Li}
\affiliation{Meta Platforms, Inc., Menlo Park, CA 94025}

\author{Naman Katyal}
\affiliation{Energy Geosciences Division, Lawrence Berkeley National Laboratory, Berkeley, CA 94720}
\DIFaddbegin \affiliation{Materials Sciences Division, Lawrence Berkeley National Laboratory, Berkeley, CA 94720}
\DIFaddend 

\author{Piotr Zarzycki}
\DIFdelbegin %DIFDELCMD < \affiliation{Materials Sciences Division, Lawrence Berkeley National Laboratory, Berkeley, CA 94720}
%DIFDELCMD < %%%
\DIFdelend \DIFaddbegin \affiliation{Energy Geosciences Division, Lawrence Berkeley National Laboratory, Berkeley, CA 94720}
\DIFaddend 

\begin{abstract}
The bond-valence model is a standard way to estimate bond strengths in
crystals, but its exponential dependence on bond length has lacked a
derivation from a specific physical interaction.  We show that this form
emerges as the leading-order limit of screened Coulomb electrostatics
and that the fitted bond-valence softness can be interpreted in terms of
an electronic screening length.  This turns bond valence from an
empirical fitting rule into a transferable descriptor of local screened
charge response across coordination environments.  The resulting theory predicts how the
bond-valence parameters should vary with ionic charge and coordination
number, and that prediction agrees with 150 fitted \DIFdelbegin \DIFdel{valences }\DIFdelend \DIFaddbegin \DIFadd{relationships }\DIFaddend from 94
cation--oxygen species, including 68 in fourfold coordination and
82 in sixfold coordination\DIFdelbegin \DIFdel{, at an abundance-weighted
coefficient of determination of 0.986}\DIFdelend \DIFaddbegin \DIFadd{.  The screening length also
sets a valence-dependent cation--oxygen contact distance that
reproduces tabulated crystal radii, with an effective oxygen size
that grows by about 0.09~\AA{} per unit of cation charge}\DIFaddend .  A comparison with
first-principles charge densities shows that the bond-valence shell
radius tracks the electronic screening cloud with coefficients of determination
of 0.9998 for ten alkali and alkaline-earth oxides and 0.967 for 21 other binary oxides
whose nearest-neighbor environments match the theory's assumptions.  The
widely used bond-valence model is thus the leading-order expression of
screened electrostatics in ionic solids.
\end{abstract}

\maketitle

Pauling's observation that the strength of a bond reaching a coordinating
anion approximates the cation charge divided by the coordination number
\cite{Pauling1929} has been an organizing principle of crystal chemistry
for nearly a century.  Brown and Altermatt \cite{Brown1985} cast this idea
as the bond-valence model: the empirical bond strength between a cation $i$
and anion $j$ separated by $R_{ij}$ is
\begin{equation}\label{eq:bv}
S_{ij}=\exp\!\left(\frac{R_0-R_{ij}}{B}\right),
\end{equation}
where $R_0$ and $B$ are parameters fitted to crystallographic data and the
valence-sum rule $\sum_j S_{ij}=z_i$ (formal charge) is enforced.  The
model is used routinely for structure validation
\DIFdelbegin \DIFdel{\mbox{%DIFAUXCMD
\cite{Brown2009,Brown2016,Gagn2015}}\hskip0pt%DIFAUXCMD
}\DIFdelend \DIFaddbegin \DIFadd{\mbox{%DIFAUXCMD
\cite{Brown2009,Brown2016} }\hskip0pt%DIFAUXCMD
and increasingly as a structural descriptor
in machine-learning models \mbox{%DIFAUXCMD
\cite{Li2021,Zhang2023,Miller2023}}\hskip0pt%DIFAUXCMD
}\DIFaddend .

Several lines of work have sought a physical foundation for the
bond-valence exponential.  Preiser \emph{et al.}\ \cite{Preiser1999}
\DIFdelbegin \DIFdel{gave
a clear Gauss's-law account, interpreting bond valence as localized
electrostatic flux , and }\DIFdelend \DIFaddbegin \DIFadd{showed that bond valences closely
track the electrostatic flux linking each cation--anion pair,
recasting the model as localized electrostatics constrained by
Gauss's law. }\DIFaddend Brown developed this \DIFdelbegin \DIFdel{picture further within a broader
electrostatic framework \mbox{%DIFAUXCMD
\cite{Brown2009,Brown2016}}\hskip0pt%DIFAUXCMD
. However, in those
}\DIFdelend \DIFaddbegin \DIFadd{flux picture into a broader
framework in which bonds act as localized capacitors
\mbox{%DIFAUXCMD
\cite{Brown2009,Brown2016}}\hskip0pt%DIFAUXCMD
.  In these }\DIFaddend treatments the exponential
distance dependence is postulated rather than derived from a
specific potential, \DIFdelbegin \DIFdel{leaving }\DIFdelend \DIFaddbegin \DIFadd{so the softness }\DIFaddend $B$ \DIFdelbegin \DIFdel{as an empirical fit
parameter with length
units but no associated physical grounding. Adams related the softness }\DIFdelend \DIFaddbegin \DIFadd{remains an empirical length
with no associated interaction. In practice the near-degeneracy of jointly
fitted $(R_0,B)$ is well documented \mbox{%DIFAUXCMD
\cite{Gagn2015}}\hskip0pt%DIFAUXCMD
, and it has been suppressed
by fixing $B\simeq0.37$~\AA\ \mbox{%DIFAUXCMD
\cite{Brown1985}}\hskip0pt%DIFAUXCMD
. Adams and co-workers connected }\DIFaddend $B$
to Morse-potential curvature and \DIFdelbegin \DIFdel{motivated }\DIFdelend \DIFaddbegin \DIFadd{showed that transferable parameter sets require
}\DIFaddend non-universal $B$ \DIFdelbegin \DIFdel{values \mbox{%DIFAUXCMD
\cite{Adams2001,Chen2017}}\hskip0pt%DIFAUXCMD
, but does not connect $B$ to a screening length.
Modern first-principles work can compute screened local interactions
from electronic structure \mbox{%DIFAUXCMD
\cite{Aryasetiawan2004,Shih2012}}\hskip0pt%DIFAUXCMD
, but does not derive the bond-valence form itself.  Here }\DIFdelend \DIFaddbegin \DIFadd{\mbox{%DIFAUXCMD
\cite{Adams2001,Chen2017}}\hskip0pt%DIFAUXCMD
.
}

\DIFadd{The $(R_0,B)$ correlation itself has not been treated as physical content of
the theory, but here }\DIFaddend we show that a Yukawa expansion obtains
the exponential as the leading term of a specific screened
potential, ties $B$ to a measurable screening length, predicts the
$B$--$R_0$ correlation with no fitted constants, and connects the
fitted parameters to electronic structure.  \DIFaddbegin \DIFadd{In what follows,
we extend Brown's postulate that screened flux partitions among coordinating anions,
from which isotropic Yukawa weight and its linearization are the two controlled
approximations. We show that partition-thermodynamic relations emerge as exact
identities of Eq.~\eqref{eq:bv}. Modern first-principles methods that compute
screened local Coulomb interactions directly from electronic structure
\mbox{%DIFAUXCMD
\cite{Aryasetiawan2004,Shih2012} }\hskip0pt%DIFAUXCMD
are shown to correlate with the Yukawa screening
length, establishing a physical foundation for the bond-valence model.
}\DIFaddend 

\DIFdelbegin \emph{\DIFdel{Screened-flux derivation.}} %DIFAUXCMD
\DIFdelend \DIFaddbegin \runinsec{\DIFadd{Screened-flux derivation.}} \DIFaddend A cation with charge $Q_i=z_i e$ emits
electrostatic flux $Q_i/\epsilon_0$.  In a \DIFdelbegin \DIFdel{Yukawa field with screening
length $\lambda$, the }\DIFdelend \DIFaddbegin \DIFadd{polarizable medium the
intervening electrons screen this flux, and the standard static
description of screening replaces the Coulomb potential by the Yukawa
form $\phi(r)\propto e^{-r/\lambda}/r$, with $\lambda$ the screening
length \mbox{%DIFAUXCMD
\cite{Debye1923,Thomas1927,Fermi1928,Yukawa1935}}\hskip0pt%DIFAUXCMD
.  Applying
Gauss's law to this potential, the flux threading a sphere of radius
$R$ is $(Q_i/\epsilon_0)\,e^{-R/\lambda}(1+R/\lambda)$, so the }\DIFaddend signed
flux amplitude captured by bond $i$--$j$ \DIFdelbegin \DIFdel{is
$\widetilde\Phi_{ij}\propto\Omega_{ij}\,e^{-R_{ij}/\lambda}(1+R_{ij}/\lambda)$, where $\Omega_{ij}$ is the solid-angle fraction assigned to neighbor $j$}\DIFdelend \DIFaddbegin \DIFadd{at length $R_{ij}$ is
$\widetilde\Phi_{ij}\propto e^{-R_{ij}/\lambda}(1+R_{ij}/\lambda)$.
Preiser }\emph{\DIFadd{et al.}}\DIFadd{\ \mbox{%DIFAUXCMD
\cite{Preiser1999} }\hskip0pt%DIFAUXCMD
partitioned flux according to
Gauss's law such that each bond carries a valence proportional to the
electrostatic flux it subtends.  We retain the same correspondence, but for
the screened field. 
}

\DIFadd{Adopting this single-$\lambda$ Yukawa weight for every bond of the
shell is the physical ansatz of the derivation}\DIFaddend .
Bond valence is a nonnegative partition of the cation valence among its
anion neighbors\DIFdelbegin \DIFdel{.  Imposing the }\DIFdelend \DIFaddbegin \DIFadd{, and two conditions fix that partition.  The
}\DIFaddend valence-sum rule $\sum_j S_{ij}=z_i$ \DIFdelbegin \DIFdel{and
preservation of }\DIFdelend \DIFaddbegin \DIFadd{sets the total.
Screening attenuates the flux, and the sum rule renormalizes
the total back to the formal charge, so the screened field determines only the
relative apportionment among bonds.  Preserving that
apportionment means that ratios of bond valences equal ratios of the
captured }\DIFaddend screened-flux magnitudes,
\DIFdelbegin \DIFdel{$S_{ij}/S_{ik}=w_{ij}/w_{ik}$,
uniquely gives
}\begin{displaymath}
\DIFdel{S_{ij}=z_i\,\frac{w_{ij}}{\sum_k w_{ik}},
\qquad w_{ij}:=|\widetilde\Phi_{ij}|.
}\end{displaymath}%DIFAUXCMD
\DIFdelend \DIFaddbegin \begin{equation*}
\DIFadd{\frac{S_{ij}}{S_{ik}}=\frac{w_{ij}}{w_{ik}},
\qquad w_{ij}:=|\widetilde\Phi_{ij}|.
}\end{equation*}\DIFaddend 
\DIFdelbegin \DIFdel{Choose $R^*$ operationally as the Boltzmann-weighted mean bond length
$m_i:=\sum_j p_{ij}R_{ij}$ (defined fully in the partition-function
section below).  Expanding }\DIFdelend \DIFaddbegin \DIFadd{This guarantees the non-negativity that Eq.~\eqref{eq:bv} requires,
imports the distance dependence of the screened interaction into the
partition, and introduces a single screening parameter, $\lambda$, to
scale the distance dependence.  Together the two conditions uniquely give
}\begin{equation}\DIFadd{\label{eq:partition}
S_{ij}=z_i\,\frac{w_{ij}}{\sum_k w_{ik}}.
}\end{equation}
\DIFadd{To linearize, we expand }\DIFaddend $\ln w_{ij}$ about \DIFdelbegin \DIFdel{$R^*$ gives
}\begin{displaymath}\DIFdel{%DIFDELCMD < \label{eq:expansion}%%%
\ln w_{ij}= C_i - \frac{R^*}{\lambda(\lambda+R^*)}\,(R_{ij}-R^*)+\delta_{ij},
}\end{displaymath}%DIFAUXCMD
\DIFdelend \DIFaddbegin \DIFadd{a reference bond length $R'$
chosen within the coordination shell, giving
}\begin{equation}\DIFadd{\label{eq:expansion}
\ln w_{ij}= C_i - \frac{R'}{\lambda(\lambda+R')}\,(R_{ij}-R')+\delta_{ij},
}\end{equation}\DIFaddend 
where $C_i$ is independent of bond index $j$ and
\DIFdelbegin \DIFdel{$\delta_{ij}=O[(R_{ij}-R^*)^2/\{\lambda(\lambda+R^*)\}]$}\DIFdelend \DIFaddbegin \DIFadd{$\delta_{ij}=O[(R_{ij}-R')^2/(\lambda+R')^2]$}\DIFaddend .
Exponentiating the linear
term and renormalizing therefore reproduces Eq.~\eqref{eq:bv} at
\DIFdelbegin \DIFdel{$R_{ij}=R^*$ }\DIFdelend \DIFaddbegin \DIFadd{$R_{ij}=R'$ }\DIFaddend with only quadratic log-weight error away from \DIFdelbegin \DIFdel{$R^*$}\DIFdelend \DIFaddbegin \DIFadd{$R'$}\DIFaddend .  Define
the leading-order Yukawa exponent denominator
\DIFdelbegin \begin{displaymath}\DIFdel{\label{eq:B_composite}
B = \frac{\lambda(\lambda+R^*)}{R^*}.
}\end{displaymath}%DIFAUXCMD
\DIFdelend \DIFaddbegin \begin{equation}\DIFadd{\label{eq:B_composite}
B = \frac{\lambda(\lambda+R')}{R'}.
}\end{equation}\DIFaddend 
This $B$ is the first-order exponent denominator produced by the Taylor
expansion of the Yukawa weight. \DIFdelbegin \DIFdel{On the screening branch it is positive
and coincides with the conventional bond-valence parameter; on the
anti-screening branch it becomes negative. }\DIFdelend The screening length $\lambda$ is the
primary physical variable, \DIFdelbegin \DIFdel{while $B$ is its leading-order image in the bond-valence
exponential.  On the screening branch (}\DIFdelend \DIFaddbegin \DIFadd{and the sign of $\lambda$ distinguishes
screening from anti-screening.  For }\DIFaddend $\lambda>0$ \DIFdelbegin \DIFdel{, $B>0$), inverting
Eq.~\eqref{eq:B_composite}
at the shell level gives the structure-specific screening length
$\lambda=(-R^*+\sqrt{R^{*2}+4BR^*})/2$.
At the species-level
characteristic pair defined below, we evaluate the same positive-branch
inversion on the characteristic reference shell, for which the operational choice $R^*=m_i$ gives $R^*=R_0^*$, so the tabulated
species constant is $\lambda^*=(-R_0^*+\sqrt{R_0^{*2}+4B^*R_0^*})/2$ for species with
$B^*>0$.  }\DIFdelend \DIFaddbegin \DIFadd{the Yukawa weight
decays with bond length, so flux reaching the more distant anions is
attenuated and shorter bonds capture the larger share of the valence.
For $-R'<\lambda<0$ the weight instead grows with bond length, and
longer bonds carry the larger share.  Fits return this anti-screening
sign for 16\% of the coordination shells in the corpus analyzed below,
but roughly half of those are nearly degenerate shells whose fitted
$B$ is barely constrained (Supplemental Material \mbox{%DIFAUXCMD
\cite{suppmat}}\hskip0pt%DIFAUXCMD
).  The
chemically systematic cases are square-planar and linear late
transition metals (Pd$^{2+}$, Au$^{3+}$, linear O--Cu--O) and large,
soft cations in high, irregular coordination, where directional
bonding or ambiguous coordination produces anisotropic flux capture
outside the isotropic Yukawa assumption.
}\DIFaddend 

\DIFdelbegin \emph{\DIFdel{Parameter-free slope.}} %DIFAUXCMD
\DIFdelend \DIFaddbegin \runinsec{\DIFadd{Parameter-free slope.}} \DIFaddend In the uniform-bond limit, where
$R_{ij}=\bar R$ for all $j$, the valence-sum rule gives
$R_0=\bar R+B\ln(z/n)$.  Differentiating:
\begin{equation}\label{eq:pole}
\beta := \frac{dB}{dR_0} = \frac{1}{\ln(z/n)}.
\end{equation}
\DIFdelbegin \DIFdel{This is parameter-free: the }\DIFdelend \DIFaddbegin \DIFadd{The }\DIFaddend slope of the $B$--$R_0$ correlation at fixed
coordination number $n$ is determined entirely by the charge-to-coordination
ratio $z/n$, with a pole at $z=n$\DIFdelbegin \DIFdel{marking the screening/anti-screening crossover}\DIFdelend \DIFaddbegin \DIFadd{, with no fitting parameters}\DIFaddend .
Figure~\ref{fig:pole} shows that Eq.~\eqref{eq:pole} reproduces the
fitted slopes for 94 cation--oxygen species at $n=4$ and $n=6$ (68 at
$n=4$, 82 at $n=6$\DIFdelbegin \DIFdel{; }\DIFdelend \DIFaddbegin \DIFadd{, }\DIFaddend 150 slopes in total\DIFdelbegin \DIFdel{; weighted $R^2=0.986$, }\DIFdelend \DIFaddbegin \DIFadd{).  The weighted $R^2$ is
$0.986$, with }\DIFaddend mean absolute error $0.15$ \DIFdelbegin \DIFdel{; }\DIFdelend \DIFaddbegin \DIFadd{under the }\DIFaddend weighting scheme
defined in the Supplemental Material \cite{suppmat}\DIFaddbegin \DIFadd{.  Each fitted
slope is the regression of one species' per-structure $(R_0,B)$ fits at fixed $n$,
constructed as described in the empirical validation below.  The
slopes carry bootstrap standard errors, with median relative
standard error $0.2\%$ across the 150 slopes and 90th percentile
$4\%$, so the scatter at fixed $(z,n)$---most visible near
$z=3$---reflects real differences among species rather than fit
uncertainty (per-species slopes with error bars in the Supplemental
Material \mbox{%DIFAUXCMD
\cite{suppmat}}\hskip0pt%DIFAUXCMD
}\DIFaddend ).

\DIFdelbegin \emph{\DIFdel{Screening collapse at $z=n$.}} %DIFAUXCMD
\DIFdel{The pole has direct physical content.
}\DIFdelend \DIFaddbegin \runinsec{\DIFadd{Origin of the near-pole spread.}}  \DIFadd{The spread in
Fig.~\ref{fig:pole} concentrates in the charge classes nearest each
pole, and the pole form itself sets that pattern.  At leading order
$1/\beta=\ln(z/n)$, and the second-order terms $\delta_{ij}$ of
Eq.~\eqref{eq:expansion} perturb this inverse slope, so a
species-level perturbation $u$ shifts the fitted slope by
$-\beta^{2}u$.  Fixed chemical scatter is thus amplified
quadratically toward the pole, and across the charge classes of
Fig.~\ref{fig:pole} the per-class spread indeed grows as
$|\beta|^{1.8}$.  The $z=3$ class at $n=4$ is the clearest case.
The same cations at $n=6$ match Eq.~\eqref{eq:pole} to a median
$|\Delta\beta|=0.03$, so chemical diversity alone is not the
origin.  The largest deviations belong to species for which
fourfold oxygen coordination is rare, chiefly the trivalent
lanthanides and Y (median 17 structures per fit), whereas the
well-populated trivalents Ga$^{3+}$ (209 structures), B$^{3+}$
(214), and Fe$^{3+}$ (160) lie on the curve.  The deviant slopes
are almost all shallower than predicted, consistent with
$\delta_{ij}$ attenuating the leading-order link in sparse,
atypical shells, and the structure-count weighting downweights
such shells on physical grounds.
}

\runinsec{\DIFadd{Screening collapse at $z=n$.}} \DIFadd{What degenerates at $z=n$ is the
leading-order link between the fitted parameters.
The $B$ dependence of $R_0$ vanishes as
$z\to n$, so the intercept ceases to constrain the softness.  $B$ is
then fixed only by the second-order spread terms $\delta_{ij}$ in
Eq.~\eqref{eq:expansion}.  The branch change across the pole is
carried by the sign reversal of the slope $\beta=1/\ln(z/n)$, negative
for $z<n$ and positive for $z>n$. }\DIFaddend At $z=n$ \DIFdelbegin \DIFdel{every bond carries exactly one valence unit
(}\DIFdelend \DIFaddbegin \DIFadd{the sum rule fixes only the
mean bond valence at unity, but it does not force }\DIFaddend $S_{ij}=1$ \DIFdelbegin \DIFdel{), so $\exp[(R_0-R_{ij})/B]=1$ for all~$j$.For a symmetric shell with $R_{ij}=\bar R$ this is
trivially satisfied at }\DIFdelend \DIFaddbegin \DIFadd{bond
by bond.  For a truly degenerate shell---all bond lengths equal by
symmetry, as for Si$^{4+}$ in zircon (ZrSiO$_4$), whose four Si--O
bonds are symmetry-equivalent---Eq.~\eqref{eq:bv}
assigns every bond the same valence, and a single structure supplies
only the uniform-limit relation $R_0=\bar R+B\ln(z/n)$.  $B$ is
unconstrained, and at $z=n$ the relation collapses to $R_0=\bar R$ at
}\DIFaddend any $B$\DIFdelbegin \DIFdel{with $R_0=\bar R$---the $z=n$ limit is degenerate}\DIFdelend .  For a \DIFdelbegin \DIFdel{real
coordination shell with unequal bond lengths
}\DIFdelend \DIFaddbegin \DIFadd{nearly degenerate shell the bond lengths
differ, so the $S_{ij}$ must differ and the sum rule is
satisfied at finite $B$.  $\alpha$-quartz is a common
example, with two experimental Si--O distances differing by about
$0.01$~\AA{} that finite $B$ resolves.  Both cases are common.  Of
the 7,565 shells with $z=n$ in our corpus, 959 are
degenerate within $10^{-3}$~\AA{} and 2}\DIFaddend ,\DIFdelbegin \DIFdel{however, the only solution
is
$B\to\infty$, so that every finite $(R_0-R_{ij})/B\to 0$.  On the screening branch,Eq.
~\eqref{eq:B_composite} gives
$\lambda\approx\sqrt{BR^*}\to\infty$: the screening length diverges.  On the $z>n$ side of the pole}\DIFdelend \DIFaddbegin \DIFadd{960 more have bond-length
spreads within $0.05$~\AA{}.
}

\begin{figure}[tb]
\centering
\includegraphics[width=\columnwidth]{prm_oxygen_beta_vs_charge_box_cn4_cn6.png}
\caption{\DIFaddFL{Fitted slope $\beta$ vs.\ formal charge $z$ for 94
cation--oxygen species at (a) $n=4$ (68 cations) and (b) $n=6$ (82
cations). 56 cations contribute at both coordinations, giving
150 fitted slopes in total. Boxes span the interquartile range of the
per-species slopes at each $z$, black points mark the median,
whiskers extend to the most extreme slopes within $1.5\times$ the
interquartile range, and open circles mark slopes beyond the whiskers.
Dashed curves: $\beta=1/\ln(z/n)$.  The
pole at $z=n$ separates the negative-slope ($z<n$) and positive-slope
($z>n$) branches of Eq.~\eqref{eq:pole}. Species at $z=n$ are omitted.
Per-species slopes are resolved in the Supplemental
Material \mbox{%DIFAUXCMD
\cite{suppmat}}\hskip0pt%DIFAUXCMD
.}}
\label{fig:pole}
\end{figure}

\begin{figure*}[t]
\centering
\includegraphics[width=\textwidth]{thomas_fermi_reff_prl.png}
\caption{\DIFaddFL{First-principles screening centroid $r_{\mathrm{eff}}$ vs.\ Gibbs shell
centroid $m_i$.  (a) Ten $s$-block ground-state binary oxides.  The solid
line is the shared fit and the dashed line is $r_{\mathrm{eff}}=m_i$.  (b)
Twenty-five non-$s$-block ground-state binary oxides, color-coded by
block.  Open symbols mark the four directional-bonding outliers
excluded from the regression (B$_2$O$_3$, P$_2$O$_5$, As$_2$O$_3$,
Au$_2$O$_3$).}}
\label{fig:centroid}
\end{figure*}

\runinsec{\DIFadd{Characteristic intersections.}} \DIFadd{Species that sit at or near
$z=n$, e.g., B$^{3+}$ in trigonal coordination (B$_2$O$_3$, band gap
6~eV), C$^{4+}$ at $n=4$, and Si$^{4+}$ in tetrahedral coordination
(SiO$_2$}\DIFaddend , \DIFaddbegin \DIFadd{9~eV), lie in the wide-gap covalent limit where weak
screening is physically plausible.  For these species the pole marks
the breakdown of the first-order link between $R_0$ and }\DIFaddend $B$\DIFdelbegin \DIFdel{changes sign ($B<0$) }\DIFdelend \DIFaddbegin \DIFadd{, so the
$B$--$R_0$ line fitted at the pole coordination cannot by itself
constrain the softness.  The data, however, resolve each cation at
more than one coordination number, and the $n$-resolved $B$--$R_0$
lines of a given cation (slopes $\beta^{(n)}$) are not independent.
Plotted together, they intersect at a nearly common point
(Supplemental Material \mbox{%DIFAUXCMD
\cite{suppmat}}\hskip0pt%DIFAUXCMD
).  Because the slopes are fixed
by $z/n$, the intersection is set entirely by the line intercepts.
Writing each $n$-resolved line as $B=\beta^{(n)}R_0+\beta_0^{(n)}$,
\mbox{%DIFAUXCMD
\cite{Li2025} }\hskip0pt%DIFAUXCMD
the uniform-limit relation $R_0=\bar R+B\ln(z/n)$ gives
}\begin{equation}\DIFadd{\label{eq:intercept}
\beta_0^{(n)}=-\frac{\bar R^{(n)}}{\ln(z/n)},
}\end{equation}
\DIFadd{where $\bar R^{(n)}$ is the characteristic shell radius of the
species at coordination $n$.  Each line crosses $B=0$ at the shell
radius itself, so the slope carries the charge-to-coordination ratio
while the intercept records where the shell sits.  The intersection
of two coordination branches $n_1<n_2$ is therefore fixed by the
dilation of the shell with coordination number,
}\begin{equation}\DIFadd{\label{eq:bstar_dilation}
B^*=\frac{\bar R^{(n_2)}-\bar R^{(n_1)}}{\ln(n_2/n_1)},
}\end{equation}
\DIFadd{with $R_0^*=\bar R^{(n_1)}+B^*\ln(z/n_1)$ the common radius at which
the branches agree.  Bond lengths grow with coordination number, so
Eq.~\eqref{eq:bstar_dilation} gives $B^*>0$ naturally on the
screening branch, and the characteristic softness acquires a direct
geometric meaning.  It is the logarithmic rate at which the
coordination shell dilates with $n$}\DIFaddend .
\DIFdelbegin \DIFdel{Its magnitude
still diverges, so the branch change is carried directly by $\mathrm{sgn}\,B$
or,
equivalently, by the sign reversal of $\beta=1/\ln(z/n)$.  Approaching the
pole from either
side ($z=n\mp\varepsilon$, $\varepsilon\to 0$), the divergence rate is
governed by the
species characteristic }\DIFdelend \DIFaddbegin 

\DIFadd{These intersections define the characteristic pair $(R_0^*,B^*)$,
species-level constants fixed by the family of lines rather than by
any single member, and therefore well defined even when one line
steepens toward the pole.  Two lines on the same branch differ in
slope by $\ln(n_2/n_1)/[\ln(z/n_1)\ln(z/n_2)]$, which is small
whenever $z$ is well separated from both coordination numbers, so
such lines can appear nearly coincident.  Their intersection
nonetheless remains constrained exactly to the extent that the data
resolve the shell dilation in Eq.~\eqref{eq:bstar_dilation}.
Operationally, $(R_0^*,B^*)$ minimizes the weighted variance of the
coordination-resolved lines, with each line weighted by the number
of structures retained in its final fit, and the pooled construction
reports the residual spread $\sigma_B$ as its conditioning
diagnostic (Supplemental Material \mbox{%DIFAUXCMD
\cite{suppmat}}\hskip0pt%DIFAUXCMD
).  The
screening-branch $B^*$ values are not universal but cluster near the
conventional $B=0.37$~\AA{} \mbox{%DIFAUXCMD
\cite{Brown1985}}\hskip0pt%DIFAUXCMD
, and their comparison
with published softness parameterizations
\mbox{%DIFAUXCMD
\cite{Adams2001,Chen2017,Gagn2015} }\hskip0pt%DIFAUXCMD
is given in the Supplemental
Material \mbox{%DIFAUXCMD
\cite{suppmat}}\hskip0pt%DIFAUXCMD
.
}

\DIFadd{On the screening branch ($\lambda>0$, $B>0$), inverting
Eq.~\eqref{eq:B_composite} at the shell level gives the
structure-specific screening length
$\lambda=(-R'+\sqrt{R'^2+4BR'})/2$.  Evaluating the same
positive-branch inversion at the characteristic point, where the
operational choice of $R'$ equals the characteristic }\DIFaddend intercept
$R_0^*$\DIFdelbegin \DIFdel{(the $n$-resolved-line convergence point introduced below):
}\begin{displaymath}
\DIFdel{|B| \sim \frac{n\,|R_0-R_0^*|}{\varepsilon},
\qquad
|\lambda| \sim \sqrt{\frac{R^*\,n\,|R_0-R_0^*|}{\varepsilon}}.
}\end{displaymath}%DIFAUXCMD
\DIFdel{The screening length diverges as $\varepsilon^{-1/2}$.  In }\DIFdelend \DIFaddbegin \DIFadd{, gives the tabulated species constant
$\lambda^*=(-R_0^*+\sqrt{R_0^{*2}+4B^*R_0^*})/2$.  The root is real
for $B^*\ge -R_0^*/4$.  It is positive on the screening branch and
negative on the anti-screening branch.  The pole itself has a
screening interpretation.  The species at $z=n$ are wide-gap
insulators (B$_2$O$_3$ and SiO$_2$ above), and in }\DIFaddend conventional
screening theory \DIFdelbegin \DIFdel{, metals have }\DIFdelend \DIFaddbegin \DIFadd{gapped systems screen weakly because low-energy
charge response is suppressed, in contrast to the }\DIFaddend finite
Thomas--Fermi screening set by a \DIFdelbegin \DIFdel{nonzero }\DIFdelend \DIFaddbegin \DIFadd{metal's }\DIFaddend density of states at the
Fermi level \DIFdelbegin \DIFdel{, whereas gapped or
correlated insulators suppress low-energy charge response and therefore
screen more weakly \mbox{%DIFAUXCMD
\cite{Thomas1927,Fermi1928,Debye1923,AshcroftMermin1976}}\hskip0pt%DIFAUXCMD
.
}\DIFdelend \DIFaddbegin \DIFadd{\mbox{%DIFAUXCMD
\cite{Thomas1927,Fermi1928,Debye1923,AshcroftMermin1976}}\hskip0pt%DIFAUXCMD
.
The fitted family behaves the same way.  Approaching the pole
($z=n\mp\varepsilon$, $\varepsilon\to 0$), sustaining a finite
intercept offset $R_0-R_0^*$ requires the inferred screening length
to grow as $|\lambda|\sim\varepsilon^{-1/2}$, so the species where
weak screening is physically expected are exactly those for which
the model demands an unbounded screening length.  }\DIFaddend Our coarse-grained
$\lambda$ is not identical to those microscopic lengths, but \DIFdelbegin \DIFdel{its pole at $z=n$ is consistent with the same weak-screening
limit: the bond-valence partition becomes insensitive to finite
bond-length differences and }\DIFdelend the \DIFdelbegin \DIFdel{fitted $B$--$R_0$ parameters cease to
resolve the shell.  Species that sit at or near $z=n$---Si$^{4+}$ in tetrahedral
coordination (SiO$_2$,
band gap 9~eV), B}\DIFdelend \DIFaddbegin \DIFadd{two
weak-screening limits agree.
}

\DIFadd{The $\lambda^*$ values determined from the $B^*$ are tabulated in
the Supplemental Material \mbox{%DIFAUXCMD
\cite{suppmat} }\hskip0pt%DIFAUXCMD
for all 103 species,
positive for the 101 screening-branch species and negative for the
two anti-screening species (Mo}\DIFaddend $^{3+}$\DIFdelbegin \DIFdel{in trigonal
coordination (B$_2$O$_3$, 6~eV), and C$^{4+}$ at $n=4$---lie in the
wide-gap covalent limit where such weak screening is physically
plausible.  }\DIFdelend \DIFaddbegin \DIFadd{, Sb$^{3+}$), whose departures
from the first-order Yukawa regime are diagnosed there.  The two
roots of the species-level screening length also have physical
meaning.  The anti-screening alias of $\lambda^*$ has magnitude
$|\lambda^*_-|=R_0^*+\lambda^*$.  This is the characteristic shell
radius extended by one screening length, the outer reach of the
screened bond.  Compared against Shannon crystal radii
\mbox{%DIFAUXCMD
\cite{Shannon1976}}\hskip0pt%DIFAUXCMD
, $|\lambda^*_-|$ behaves as a valence-dependent
contact distance (Supplemental Material \mbox{%DIFAUXCMD
\cite{suppmat}}\hskip0pt%DIFAUXCMD
).  Here
$r_{\mathrm{cr}}$ denotes the cation's Shannon crystal radius
averaged over its observed coordination numbers, the geometric
baseline for the cation--oxygen contact.  The comparison is
empirical, with no assumption about the physics that sets the
crystallographic radii.  Within each formal-charge class the slope
of $|\lambda^*_-|$ against $r_{\mathrm{cr}}$ is fixed at one and a
single offset $c(z)$ is fitted,
}\begin{equation*}
\DIFadd{|\lambda^*_-|=r_{\mathrm{cr}}+c(z),
}\end{equation*}
\DIFadd{so $c(z)$ measures how far the outer reach of the screened bond
extends beyond the cation radius.  These one-parameter fits describe
each class well, and the offsets grow linearly with formal charge,
$c(z)\approx 1.05+0.09\,z$~\AA.  The two terms have direct physical
readings.  The constant $1.05$~\AA{} sits at the scale of the oxygen
crystal radius, so at low charge $r_{\mathrm{cr}}+c(z)$ is the
cation radius plus an oxygen radius, the geometric cation--oxygen
contact distance.  The increment shows the screened bond reaching
about $0.09$~\AA{} beyond that contact per unit of formal charge, so
$c(z)$ acts as an effective oxygen size.  The effective oxygen size
seen by the bond-valence partition is therefore not a fixed sphere.
It grows with cation valence, which is the statement that the
bond-valence shell follows the electronic screening cloud rather
than a hard ionic contact.  Comparison to first-principles electron
densities tests this directly, but first requires a shell
radius that the bond-valence partition itself defines.
}\DIFaddend 

\DIFdelbegin %DIFDELCMD < \begin{figure}[tb]
%DIFDELCMD < \centering
%DIFDELCMD < \includegraphics[width=\columnwidth]{prl_oxygen_beta_vs_charge_cn4_cn6.png}
%DIFDELCMD < %%%
%DIFDELCMD < \caption{%
{%DIFAUXCMD
\DIFdelFL{Fitted slope $\beta$ vs.\ formal charge $z$ for 94
cation--oxygen species at $n=4$ (68 cations, left) and $n=6$ (82
cations, right); 56 cations contribute at both coordinations, giving
150 fitted slopes in total.  Dashed curves: $\beta=1/\ln(z/n)$.  The
pole at $z=n$ separates the screening ($z<n$) and anti-screening
($z>n$) branches; species at $z=n$ are omitted because the $B$--$R_0$
fit is degenerate at the pole (see Supplemental Material
\mbox{%DIFAUXCMD
\cite{suppmat}}\hskip0pt%DIFAUXCMD
).}}
%DIFAUXCMD
%DIFDELCMD < \label{fig:pole}
%DIFDELCMD < \end{figure}
%DIFDELCMD < 

%DIFDELCMD < \begin{figure*}[t]
%DIFDELCMD < \centering
%DIFDELCMD < \includegraphics[width=\textwidth]{thomas_fermi_reff_prl.png}
%DIFDELCMD < %%%
%DIFDELCMD < \caption{%
{%DIFAUXCMD
\DIFdelFL{First-principles screening centroid $r_{\mathrm{eff}}$ vs.\ Boltzmann shell
centroid $m_i$.  (a) Ten $s$-block ground-state binary oxides.  Solid
line: shared fit; dashed line: $r_{\mathrm{eff}}=m_i$.  (b)
Twenty-five non-$s$-block ground-state binary oxides, color-coded by
block.  Open symbols mark the four directional-bonding outliers
excluded from the regression (B$_2$O$_3$, P$_2$O$_5$, As$_2$O$_3$,
Au$_2$O$_3$).}}
%DIFAUXCMD
%DIFDELCMD < \label{fig:centroid}
%DIFDELCMD < \end{figure*}
%DIFDELCMD < 

%DIFDELCMD < %%%
\emph{\DIFdel{Partition thermodynamics.}} %DIFAUXCMD
\DIFdel{Because }\DIFdelend \DIFaddbegin \runinsec{\DIFadd{Partition thermodynamics.}} \DIFadd{Reading
Eq.~\eqref{eq:partition} as a statistical ensemble supplies a characteristic
shell radius and also fixes the reference length $R'$ left open above.
These are the direct consequence of Eqs.~\eqref{eq:bv} and
\eqref{eq:partition}.  We introduce thermodynamic language as interpretation,
and develop the structure in three
steps.  First, normalize the bond valences within the shell:
}\begin{equation}\DIFadd{\label{eq:pweights}
p_{ij}:=\frac{S_{ij}}{\sum_k S_{ik}}
=\frac{w_{ij}}{\sum_k w_{ik}}
=\frac{e^{-R_{ij}/B}}{Z_i},
}\end{equation}
\DIFadd{where $Z_i:=\sum_j e^{-R_{ij}/B}$.  The middle form follows from
Eq.~\eqref{eq:partition}, whose shell total is $z_i$, so $p_{ij}$ is
the normalized screened-flux weight of bond $j$.  The last form
follows from Eq.~\eqref{eq:bv}, in which }\DIFaddend $R_0/B$ cancels in \DIFdelbegin \DIFdel{$S_{ij}/\sum_k S_{ik}$, the
normalized bond weights
$p_{ij}=S_{ij}/\sum_k S_{ik}=e^{-R_{ij}/B}/Z_i$ define a }\DIFdelend \DIFaddbegin \DIFadd{the
ratio.  The middle and last forms coincide to the first order of
Eq.~\eqref{eq:expansion}.  Equation~\eqref{eq:pweights} is a Gibbs distribution over the
coordination shell.  Bond length plays the role of energy, $Z_i$ is
the }\DIFaddend canonical partition function\DIFdelbegin \DIFdel{with $B>0$ as }\DIFdelend \DIFaddbegin \DIFadd{, and $B$ is }\DIFaddend the temperature-like
scale\DIFdelbegin \DIFdel{on the
screening branch. The Boltzmann
mean $m_i:=\sum_j p_{ij}R_{ij}$ and the Shannon entropy
$H_i:=-\sum_j p_{ij}\ln p_{ij}$ yield the exact identity
}\begin{displaymath}
\DIFdel{R_0 = m_i + B(\ln q_i - H_i),
}\end{displaymath}%DIFAUXCMD
\DIFdelend \DIFaddbegin \DIFadd{. Second, define two coordination shell properties,
}\begin{equation}\DIFadd{\label{eq:mean_entropy}
m_i:=\sum_j p_{ij}R_{ij},
\qquad
H_i:=-\sum_j p_{ij}\ln p_{ij},
}\end{equation}
\DIFadd{the Gibbs mean---the valence-weighted centroid of the
coordination shell---and the Shannon entropy of the bond-weight
distribution.  Third, invert Eq.~\eqref{eq:bv} bond by bond,
$R_0=R_{ij}+B\ln S_{ij}$, and average over the weights $p_{ij}$.
With $\ln S_{ij}=\ln q_i+\ln p_{ij}$, the average gives the exact
identity
}\begin{equation}\DIFadd{\label{eq:R0_identity}
R_0 = m_i + B(\ln q_i - H_i),
}\end{equation}\DIFaddend 
where $q_i:=\sum_j S_{ij}$\DIFaddbegin \DIFadd{, equal to $z_i$ when the valence-sum rule
is enforced exactly}\DIFaddend .  Here $H_i$ \DIFdelbegin \DIFdel{is a purely geometric quantity:
it }\DIFdelend measures how uniformly the
bond-valence weight is distributed across the coordination shell\DIFdelbegin \DIFdel{, reaching }\DIFdelend \DIFaddbegin \DIFadd{.  It
reaches }\DIFaddend $\ln n$ for a symmetric shell and \DIFdelbegin \DIFdel{falling }\DIFdelend \DIFaddbegin \DIFadd{falls }\DIFaddend below it as distortion
concentrates weight on the \DIFdelbegin \DIFdel{shorter }\DIFdelend \DIFaddbegin \DIFadd{extremal }\DIFaddend bonds.
Thus\DIFaddbegin \DIFadd{, }\DIFaddend $R_0$ is \DIFdelbegin \DIFdel{not the mean bond length but }\DIFdelend an entropy-corrected weighted
mean\DIFdelbegin \DIFdel{, and the operational choice $R^*=m_i$ }\DIFdelend \DIFaddbegin \DIFadd{. It reduces to the unit-valence limit $R_0=R_{ij}$ that Brown
}\emph{\DIFadd{et al.}}\DIFadd{\ \mbox{%DIFAUXCMD
\cite{Brown2016} }\hskip0pt%DIFAUXCMD
derived for a single bond with $S_{ij}=1$.
The Gibbs mean also fixes the reference length left open in
the screened-flux derivation. We take $R'=m_i$, which centers the
log-weight expansion of Eq.~\eqref{eq:expansion} on the weighted
shell centroid that the partition function singles out and makes the
operational content of $R'$ }\DIFaddend in Eq.~\eqref{eq:B_composite} \DIFdelbegin \DIFdel{is made }\DIFdelend explicit.

\DIFdelbegin \emph{\DIFdel{Empirical validation.}} %DIFAUXCMD
\DIFdelend \DIFaddbegin \runinsec{\DIFadd{Empirical validation.}} \DIFaddend Bond-valence parameters $(R_0,B)$ were
determined for 103 cation--oxygen species using first-principles
relaxed crystal
structures from the Materials Project \cite{MaterialsProject,Horton2025}.
For each structure, bond strengths were obtained by solving the
valence-sum and loop constraints on the cation--oxygen bond graph
\cite{Li2025}, and $(R_0,B)$ were fitted \DIFaddbegin \DIFadd{jointly }\DIFaddend by linear least
squares\DIFdelbegin \DIFdel{; the
few rank-deficient bond graphs and numerically ill-conditioned fits
were treated by the procedure described in
the Supplemental Material
}\DIFdelend \DIFaddbegin \DIFadd{, not by re-optimizing $B$ at fixed $R_0$.
The coordination-resolved
lines of Fig.~\ref{fig:pole}, first introduced in
Ref.~\mbox{%DIFAUXCMD
\cite{Li2025}}\hskip0pt%DIFAUXCMD
, are then obtained by regressing $B$ on
$R_0$ across the per-structure pairs of each species at fixed $n$ (at
least five structures per cell), using an outlier-robust RANSAC
estimator refit by least squares on its inliers.  A fitted slope
$\beta^{(n)}$ is accepted only with at least four inlier structures,
slope uncertainty $\sigma_\beta\le 5$, and $R^2\ge 0.3$
}\DIFaddend \cite{suppmat,YukawaScreeningRepo}.  The resulting set spans $s$-,
\DIFdelbegin \DIFdel{$d$}\DIFdelend \DIFaddbegin \DIFadd{$p$}\DIFaddend -, \DIFdelbegin \DIFdel{$f$}\DIFdelend \DIFaddbegin \DIFadd{$d$}\DIFaddend -, and \DIFdelbegin \DIFdel{$p$}\DIFdelend \DIFaddbegin \DIFadd{$f$}\DIFaddend -block cations\DIFdelbegin \DIFdel{; Fig.  }\DIFdelend \DIFaddbegin \DIFadd{.  Figure}\DIFaddend ~\ref{fig:pole} shows the 94
species (150 fitted slopes) for which the prediction is finite at
$n=4$ and/or $n=6$.

\DIFdelbegin \DIFdel{Within each $(z,n)$ group, the slope $\beta^{(n)}$ confirms
Eq.~\eqref{eq:pole} across all families.  The $n$-resolved lines
converge to cation-specific characteristic pairs
$(R_0^*,B^*)$, defined as the intersection that minimizes
the weighted variance of the coordination-resolved $B$--$R_0$ lines,
with each line weighted by the final fit population of that
coordination shell (Supplemental Material \mbox{%DIFAUXCMD
\cite{suppmat}}\hskip0pt%DIFAUXCMD
).Of the 103 species, 101 have $B^*>0$ (screening branch);
for these the screening length
$\lambda^*=(-R_0^*+\sqrt{R_0^{*2}+4B^*R_0^*})/2$ is a tabulated
constant (Supplemental Material \mbox{%DIFAUXCMD
\cite{suppmat}}\hskip0pt%DIFAUXCMD
).  The remaining two
species (Mo$^{3+}$, Sb$^{3+}$) have $B^*<0$, fall on the anti-screening branch, and are
retained only as signed diagnostics.  Both sit outside the first-order
Yukawa regime: Sb$^{3+}$ has a stereoactive $5s^2$ lone pair driving
strong directional distortion, and Mo$^{3+}$ is constrained by
only two
coordination shells (Supplemental Material \mbox{%DIFAUXCMD
\cite{suppmat}}\hskip0pt%DIFAUXCMD
), so its
$(R_0^*,B^*)$ is provisional.  The screening-branch $B^*$ values are non-universal
(mean $0.38\pm0.15$~\AA), in qualitative agreement with Adams's softBV
parameterization \mbox{%DIFAUXCMD
\cite{Adams2001}}\hskip0pt%DIFAUXCMD
, while $\lambda^*$ provides a
single-parameter species descriptor.
}%DIFDELCMD < 

%DIFDELCMD < %%%
\emph{\DIFdel{First-principles confirmation.}} %DIFAUXCMD
\DIFdel{The connection to electronic structure is tested
by comparing the }\DIFdelend \DIFaddbegin \runinsec{\DIFadd{First-principles confirmation.}} \DIFadd{If bond valence is captured
screened flux, the shell radius resolved by the }\DIFaddend bond-valence
\DIFdelbegin \DIFdel{Boltzmann shell centroid
$m_i=\sum_j p_{ij}R_{ij}$ }\DIFdelend \DIFaddbegin \DIFadd{partition---the Gibbs shell centroid $m_i$ of
Eq.~\eqref{eq:mean_entropy}---should track the radius at which
screening charge accumulates along each bond.  We test this by
comparing $m_i$ }\DIFaddend to a first-principles screening centroid
$r_{\mathrm{eff}}$\DIFaddbegin \DIFadd{, the centroid of the excess electron density on
the oxygen side of the bond path, }\DIFaddend extracted from charge densities in
the Materials Project \cite{MaterialsProject,Horton2025}.  For each M--O
bond with $R_{ij}\le 3.5$~\AA, the total valence-electron density
$\rho(r)$ is sampled at 200 uniformly spaced points along the
nearest-image straight M--O segment.  An interior density-minimum proxy
on $0.15R_{ij}<r<0.85R_{ij}$ defines the start of the oxygen-side
window\DIFdelbegin \DIFdel{; }\DIFdelend \DIFaddbegin \DIFadd{, and }\DIFaddend $r_{\mathrm{eff}}$ is then the excess-density-weighted centroid
on $r_{\mathrm{BCP}}\le r\le 0.93R_{ij}$ with weights
$\max[\rho(r)-\rho_{\mathrm{BCP}},0]$, averaged over the valid
first-shell bonds of that oxide (see Supplemental Material for the exact
discrete construction and fallback rule).
For ten $s$-block binary oxides (Li$_2$O through BaO\DIFdelbegin \DIFdel{;
}\DIFdelend \DIFaddbegin \DIFadd{,
}\DIFaddend Fig.~\ref{fig:centroid}a),
these collapse onto a single line:
\begin{equation*}
r_{\mathrm{eff}} = 0.94\,m_i - 0.24~\text{\AA}
\qquad (R^2=0.9998).
\end{equation*}
This is not a trivial monotone size trend: the jackknife
mean absolute error of the $m_i$ regression is $0.006$~\AA, versus
$0.078$~\AA\ for a linear fit against $R_0$ and $0.095$~\AA\ against the
Shannon cation radius
\cite{Shannon1976}
(see Supplemental Material for the corresponding control plot).  The near-unit slope therefore shows
that the bond-valence partition function and the first-principles screening cloud track
the same ionic shell radius.  The same centroid correspondence extends
beyond the $s$ block.
Figure~\ref{fig:centroid}b shows the 25 non-$s$-block binary oxides for
which the charge density yields a well-defined oxide-level screening centroid
(listed in the Supplemental Material \cite{suppmat}).  The Yukawa
derivation assumes approximately isotropic first-shell screening, so
the regression is restricted to the 21 oxides whose cation environments
satisfy that condition, excluding four whose bonding is dominated by
directional, non-isotropic motifs: B$_2$O$_3$, P$_2$O$_5$,
As$_2$O$_3$, and Au$_2$O$_3$\DIFdelbegin \DIFdel{; the }\DIFdelend \DIFaddbegin \DIFadd{.  The }\DIFaddend Supplemental Material gives the fixed
inclusion rule and the exact oxide-level centroid construction used to
define $r_{\mathrm{eff}}$.  The 21 included oxides follow
\begin{equation*}
r_{\mathrm{eff}} = 0.98\,m_i - 0.30~\text{\AA}
\qquad (R^2=0.967).
\end{equation*}
The \DIFdelbegin \DIFdel{centroid correspondence therefore extends across }\DIFdelend \DIFaddbegin \DIFadd{correspondence therefore spans }\DIFaddend the $p$\DIFdelbegin \DIFdel{-}\DIFdelend , $d$\DIFdelbegin \DIFdel{-}\DIFdelend , and $f$ \DIFdelbegin \DIFdel{-block with $R^2=0.967$}\DIFdelend \DIFaddbegin \DIFadd{blocks}\DIFaddend , and
its breakdown is confined to strongly anisotropic directional-bonding
environments outside the isotropic Yukawa regime.

\DIFdelbegin \emph{\DIFdel{Discussion.}} %DIFAUXCMD
\DIFdelend \DIFaddbegin \runinsec{\DIFadd{Discussion.}} \DIFaddend The main result is that the bond-valence
\DIFdelbegin \DIFdel{exponential---the workhorse of crystal-chemistry validation for four
decades---is }\DIFdelend \DIFaddbegin \DIFadd{exponential is }\DIFaddend the leading-order term of Yukawa screened
electrostatics.  The fitted exponent denominator $B$ is not a free
parameter but the leading-order Yukawa denominator set by the
screening length and shell geometry, with its sign carrying the
screening/anti-screening branch information \DIFdelbegin \DIFdel{; the parameter correlation $\beta=1/\ln(z/n)$ is parameter-free }\DIFdelend \DIFaddbegin \DIFadd{that diagnoses the applicability
of the first order, linear assumption employed here. The empirical
correlation between $R_0$ }\DIFaddend and \DIFaddbegin \DIFadd{$B$ is determined by
$\beta=1/\ln(z/n)$ and }\DIFaddend exact in the uniform-bond limit\DIFdelbegin \DIFdel{;
and the }\DIFdelend \DIFaddbegin \DIFadd{.  The
}\DIFaddend partition-function structure \DIFaddbegin \DIFadd{for inequivalent bonds }\DIFaddend gives an
entropy-corrected identity for $R_0$ and a free energy
$F=m_i-BH_i=R_0-B\ln q_i$ in length units.

The \DIFdelbegin \DIFdel{first-principles centroid comparison confirms that this is not merely a
mathematical analogy: the electronic screening cloud and the bond-valence
weights describe the same shell geometry.  The small deviation from the
identity line is explained by the oxygen-side valence centroid shift
measured directly from the same charge-density paths, while the null-model
controls show that the collapse is substantially sharper than a generic
monotone size trend.
}%DIFDELCMD < 

%DIFDELCMD < %%%
\DIFdel{The }\DIFdelend derivation rests on two approximations: (i) the first-order
Taylor expansion, which breaks down for highly distorted shells
($d^4$/$d^9$ Jahn--Teller octahedra), and (ii) \DIFdelbegin \DIFdel{the assumption that solid-angle
fractions $\Omega_{ij}$ are approximately uniform across the
coordination shell.  Solid-angle non-uniformity contributes an additive
$\ln\Omega_{ij}$ term to }\DIFdelend \DIFaddbegin \DIFadd{isotropic flux
capture, in which the ansatz
$w_{ij}\propto e^{-R_{ij}/\lambda}(1+R_{ij}/\lambda)$ assigns each
bond a weight that depends on its length alone.  Both assumptions
fail when directional bonding concentrates flux capture into
particular solid angles, adding a bond-dependent geometric term that
the expansion of }\DIFaddend Eq.~\eqref{eq:expansion} \DIFdelbegin \DIFdel{that is absorbed }\DIFdelend \DIFaddbegin \DIFadd{cannot absorb }\DIFaddend into $C_i$\DIFdelbegin \DIFdel{for near-regular polyhedra but partially decorrelates the bond
weights from $\lambda$ alone when the shell departs from local
symmetry}\DIFdelend ,
as in stereoactive lone-pair and strongly anisotropic environments
\cite{Bosi2014,Brown2016}.  The empirical signature of this
breakdown is visible in Fig.~\ref{fig:centroid}b: the four oxides
whose ground-state structures fail the \DIFdelbegin \DIFdel{isotropic-shell
}\DIFdelend \DIFaddbegin \DIFadd{isotropic flux-capture
}\DIFaddend assumption (B$_2$O$_3$, P$_2$O$_5$, As$_2$O$_3$, Au$_2$O$_3$) are
precisely those that depart from the centroid line, identifying
ground-state directional bonding as the empirical boundary of the
first-order Yukawa regime.  Extending the first-principles centroid
comparison to the full set of transition-metal and $f$-block oxides
beyond the subset treated here will require going beyond the
Thomas--Fermi approximation to include orbital directionality and
correlation effects.

\DIFaddbegin \DIFadd{The first-principles centroid comparison confirms that this is not
merely a mathematical analogy, but that the electronic screening
cloud and the bond-valence weights describe the same shell geometry.
The offset from the identity line has a direct origin.  The Gibbs
shell centroid $m_i$ is built from internuclear M--O distances and
therefore references the oxygen nucleus, whereas $r_{\mathrm{eff}}$
tracks the oxygen-side screening charge, whose centroid is displaced
from the nucleus into the bond.  The same charge-density paths
measure this oxygen-side valence shift directly,
$\Delta_{\mathrm{O,val}}=0.39\pm0.03$~\AA{} on average, matching the
$m_i-r_{\mathrm{eff}}$ gap of the fitted lines across the shell
radii (Supplemental Material \mbox{%DIFAUXCMD
\cite{suppmat}}\hskip0pt%DIFAUXCMD
).  The null-model
controls show that the collapse is substantially sharper than a
generic monotone size trend.
}

\DIFaddend Bond valence is therefore recast as a coarse-grained observable of
local screened charge response \DIFdelbegin \DIFdel{rather than a purely empirical structural
heuristic. 
Concretely, the species-specific screening length
$\lambda^*$ is a compact tabulated descriptor that organizes the shell
thermodynamics and , through the identities above, provides a calibrated
starting point for estimating bond weights, entropy, effective coordination number, and $R_0$ in new coordination environments without
additional adjustable parameters}\DIFdelend \DIFaddbegin \DIFadd{in specific screening environments. 
Through species-level aggregation of screening across formal charge 
and coordination number, the screened bond coincides 
with the cation's Shannon crystal radius plus an effective oxygen size that grows with cation
valence}\DIFaddend .  Establishing the full predictive scope of that program beyond the families
validated here remains future work.  This connects Pauling's 1929
valence rules \cite{Pauling1929} to Yukawa's 1935 screened potential
\cite{Yukawa1935} through a quantitative bridge that has been
missing for nearly a century.

\begin{acknowledgments}
This work was supported by \DIFdelbegin \DIFdel{the }\DIFdelend the U.S.\ Department of Energy, Office of Science,
Office of Basic Energy Sciences, under Contract No.\ DE-AC02-05-CH11231
(MINES: The Science of Direct MINeral to Energy Storage Synthesis, FWP: FP00014914).
\end{acknowledgments}

\bibliography{references}

\end{document}
